\documentclass[11pt,a4paper]{article}
\usepackage{kotex}
\usepackage{fontspec}
\setmainfont{Times New Roman}
\setmainhangulfont{AppleMyungjo}
\setsanshangulfont{Apple SD Gothic Neo}
\usepackage[a4paper,top=18mm,bottom=18mm,left=20mm,right=20mm]{geometry}
\usepackage{fancyhdr}
\setlength{\headheight}{24pt}
\usepackage{enumitem}
\usepackage{mdframed}
\usepackage{booktabs}
\linespread{1.18}
\setlength{\parindent}{0pt}
\pagestyle{fancy}
\fancyhf{}
\renewcommand{\headrulewidth}{0.6pt}
\fancyhead[L]{\sffamily\bfseries 2026 고1 영어 변형 - 정답 및 해설 지문 40}
\fancyhead[R]{\sffamily 무의식적 위험 감지}
\fancyfoot[C]{\framebox[16mm][c]{\thepage}}
\newcommand{\qno}[1]{\par\vspace{10pt}{\sffamily\bfseries #1.}\ }
\newmdenv[linewidth=0.6pt,roundcorner=3pt,innertopmargin=7pt,innerbottommargin=7pt,
          innerleftmargin=9pt,innerrightmargin=9pt,skipabove=5pt,skipbelow=5pt]{pbox}
\begin{document}
\begin{center}
{\fontsize{15}{17}\selectfont\sffamily\bfseries [지문 40] 정답 및 해설}
\end{center}
\vspace{6pt}
\begin{center}
\renewcommand{\arraystretch}{1.35}
\begin{tabular}{cccc}
\toprule
문항 & 유형 & 정답 & 배점 \\
\midrule
1 & 요약문 & ③ & 2점 \\
2 & 빈칸추론 & ④ & 2점 \\
3 & 서술형 해석 & 모범답안 참조 & 2점 \\
4 & 영어 제목 & ③ & 2점 \\
\bottomrule
\end{tabular}
\end{center}

\qno{1}\textbf{요약문 - 정답: ③}
\begin{pbox}
글은 의식적 판단이 시작되기 전에 무의식적 ``danger detector''가 안전 여부를 먼저 판정한다고 설명한다. 따라서 (A)는 \textit{preconscious}가 가장 정확하다. 또한 인간은 감정적으로 먼저 판단한 뒤 그 판단을 정당화할 근거를 만들어 낸다고 했으므로 (B)는 \textit{rationalize}가 맞다.

\medskip
⑤ \textit{instinctive}는 어느 정도 그럴듯하지만, (B)의 \textit{ignore}가 본문과 반대라 오답이다. ① deliberate, ④ systematic은 의식적이고 체계적인 판단을 뜻하므로 글의 핵심과 어긋난다.
\end{pbox}

\qno{2}\textbf{빈칸추론 - 정답: ④}
\begin{pbox}
빈칸은 ``generate reasons for why it was accurate''를 대체하는 자리이다. 위험 감지기의 최초 반응이 맞았다고 이유를 만들어 붙이는 것이므로 ④ \textit{justify}가 정답이다.

\medskip
① challenge와 ③ override는 의식적으로 오류를 바로잡는 경우에 해당하지만, 본문은 우리가 흔히 그렇게 하지 않고 처음 반응을 정당화한다고 말한다. ② ignore와 ⑤ abandon도 이유를 만들어 붙인다는 흐름과 맞지 않는다.
\end{pbox}

\qno{3}\textbf{서술형 해석}
\begin{pbox}
\textbf{대상 문장}: \textit{When we see something or someone who ``feels'' dangerous, we have already launched into action internally before we have even started ``thinking.''}

\medskip
\textbf{모범답안}: 우리가 위험하다고 ``느껴지는'' 무언가나 누군가를 볼 때, 우리는 심지어 ``생각''을 시작하기도 전에 이미 내적으로 행동을 개시한 상태다.

\medskip
\textbf{채점 기준}: ``feels dangerous''를 실제 위험 판단이 아니라 위험하게 느껴지는 것으로, ``launched into action internally''를 내적으로 행동을 시작한 상태로, ``before ... thinking''을 생각하기도 전에로 옮기면 정답.
\end{pbox}
\vspace{8pt}
\qno{4}\textbf{제목 - 정답: ③}
\begin{pbox}
글은 의식적이고 합리적인 사고가 시작되기 전에 무의식적 위험 감지기가 먼저 판단하고, 인간은 그 뒤에 이유를 만들어 정당화한다고 설명한다. 따라서 감정 기반 선행 판단의 배후에 있는 무의식적 감지기를 나타낸 ③이 정답이다.
\end{pbox}

\end{document}
